\documentclass[a4paper]{article}

\title{Отчёт по лабораторной работе №2\\
<<Окислительно-восстановительные реакции>>}
\author{Можаров А.Р.}

\usepackage{cmap}
\usepackage[T2A]{fontenc}
\usepackage[utf8]{inputenc}
\usepackage[russian]{babel}
\usepackage{indentfirst}
\usepackage[hidelinks]{hyperref}
\usepackage{pgfplots}
\usepackage{amsmath}
\usepackage{amssymb}

\newcommand{\M}{\text{ М}}

\begin{document}

\maketitle

\subsection*{Оборудование}

\begin{enumerate}
    \item Штатив для пробирок
    \item Стеклянные пробирки
    \item Пипетки
    \item Груша
    \item Вытяжной шкаф (для проведения опытов с серной кислотой)
\end{enumerate}

\subsection*{Реагенты}

\begin{enumerate}
    \item $0.1\M$ раствор перманганата калия $KMnO_4$
    \item $0.05\M$ раствор сульфита натрия $Na_2SO_3$
    \item $0.5\M$ раствор иода $I_2$
    \item $0.05\M$ раствор тиосульфата натрия $Na_2S_2O_3$
    \item 5\% раствор медного купороса $CuSO_4$
    \item Металлическая стружка $Fe$ (скрепки)
    \item Концентрированный раствор серной кислоты $H_2SO_4$
    \item Раствор крахмала
\end{enumerate}

\section*{\centering Практическая часть}

Здесь и далее нумерация опытов совпадает с нумерацией практической части.
В ходе выполнения лабораторной работы были проделаны следующие опыты:

\begin{enumerate}
    \item 5 мл медного купороса с металлической стружкой
    \item 0.5 мл раствора перманганата калия с раствором сульфита натрия и
    \begin{enumerate}
        \item (щелочная среда) 1 мл раствора гидроксида натрия
        \item (кислая среда) несколько капель концентрированного раствора серной кислоты
        \item (нейтральная среда) ничего/вода
    \end{enumerate}
    \item 0.1 мл раствора иода и 2-3 каплями раствора крахмала с раствором тиосульфата натрия 
\end{enumerate}


\section*{\centering Наблюдения}

\begin{enumerate}
    \item Стружка покрылась ржавчиной,
    выделился газ (возможно в небольшом количестве) и <<осел>> на стружках,
    выпал бирюзово-зеленоватый осадок;
    \item
    \begin{enumerate}
        \item Раствор стал прозрачным, выпала тёмно-коричневая взвесь;
        \item Раствор стал прозрачным, осаждения или газообразования не замечено;
        \item Выпал тёмно-коричневый осадок (комками), раствор стал прозрачным;
    \end{enumerate}
    \item С тёмно-синего, после добавления тиосульфата, раствор стал прозрачным.
\end{enumerate}

\section*{\centering Уравнения реакций}

\begin{enumerate}
    \item %Данная реакция прошла по схеме ниже.
    \begin{gather*}
        CuSO_4 + Fe = Cu + FeSO_4 \\
        Cu^{2+} + 2e^- = Cu^0 \text{ (окислитель)} \\
        Fe^0 - 2e^- = Fe^{2+} \text{ (восстановитель)}
    \end{gather*}
    %В результате реакции образовалась медь,
    %дающая зеленоватый оттенок в окисленном состоянии.
    \item
    \begin{enumerate}
        \item
        \begin{gather*}
            Na_2SO_3 + 2KMnO_4 + 2NaOH = \\
            = K_2MnO_4 + Na_2MnO_4 + Na_2SO_4 + H_2O \\
            2Mn^{7+} + 2e^- = 2Mn^{6+} \text{ (окислитель)} \\
            S^{4+} - 2e^- = S^{6+} \text{ (восстановитель)}
        \end{gather*}
        \item \begin{gather*}
            3H2SO4 + 2KMnO4 + 5Na2SO3 = \\
            = 3H2O + K2SO4 + 5Na2SO4 + 2MnSO4 \\
            2Mn^{7+} + 10e^- = 2Mn^{2+} \text{ (окислитель)} \\
            5S^{4+} - 10e^- = 5S^{6+} \text{ (восстановитель)}
        \end{gather*} 
        \item \begin{gather*}
            2KMnO_4 + 3Na_2SO_3 + H_2O = \\
            = 2MnO_2 + 2KOH + 3Na_2SO_4 \\
            2Mn^{7+} + 6e^- = 2Mn^{4+} \text{ (окислитель)} \\
            3S^{4+} - 6e^- = 3S^{6+} \text{ (восстановитель)}
        \end{gather*}
    \end{enumerate}
    \item %В данном опыте йод в присутствии крахмала являлся индикатором реакции
    %(давал тёмно-синий окрас).
    \begin{gather*}
        I_2 + 2Na_2S_2O_3 = 2NaI + Na_2S_4O_6 \\
        2I^0 + 2e^- = 2I^{1-} \text{ (окислитель)} \\
        2S^{2+} - 2e^- = 2S^{4+} \text{ (восстановитель)}
    \end{gather*}
    %В результате реакции окрас пропал - весь йод прореагировал.
\end{enumerate}

\end{document}